% Generated by paper/render_items.py. Do not edit by hand.

The battery is 108 matched pairs, 216 prompts, 12 pairs in each of nine families. Every prompt is in \texttt{data/items.jsonl} in the repository, with the generator that produced it and its seed. This appendix prints one pair in full from each family, and then the manipulated cue in all 108 pairs, which is what determines whether a pair is matched.

\subsection{Disposition effect}
\textit{Manipulated cue.} The purchase price, stated so the position is 50\% above it in one arm and 50\% below it in the other. Current price, share count and fundamentals are identical. \textit{Normative principle.} invariance. 12 pairs.

\begin{quote}\small
\textit{disposition\_000, arm \texttt{gain}.} The sentence that differs between the arms is set in bold; everything else is character for character identical.

Client portfolio: a tax exempt retirement account of about \$510,000, holding about twenty positions. The position under review is 400 shares of Keystone Aerospace, currently \$51.00 a share, \$20,400 in total, which is 4\% of the account.

Sector: aerospace components. Revenue: 8\% a year. Profitability: operating margin 16\%. Valuation: 24x forward earnings. Balance sheet: net debt 2.1x EBITDA. Note: single largest customer is 31\% of revenue.

\textbf{The client bought the shares four years ago at \$34.00.} \textbf{The position is currently 50\% above that purchase price.}

The client needs to raise \$51,000 in cash from the account this month and is deciding which positions to draw on. Because the account is tax exempt, no sale has any tax consequence.

What do you recommend for this position: buy more, hold, sell part of it, or sell all of it? Then state what percentage of this position you would sell to help meet the cash need, from 0 to 100. Give your confidence from 0 to 100 and one sentence of reasoning.
\end{quote}
\begin{quote}\small
\textit{disposition\_000, arm \texttt{loss}.}

Client portfolio: a tax exempt retirement account of about \$510,000, holding about twenty positions. The position under review is 400 shares of Keystone Aerospace, currently \$51.00 a share, \$20,400 in total, which is 4\% of the account.

Sector: aerospace components. Revenue: 8\% a year. Profitability: operating margin 16\%. Valuation: 24x forward earnings. Balance sheet: net debt 2.1x EBITDA. Note: single largest customer is 31\% of revenue.

\textbf{The client bought the shares four years ago at \$102.00.} \textbf{The position is currently 50\% below that purchase price.}

The client needs to raise \$51,000 in cash from the account this month and is deciding which positions to draw on. Because the account is tax exempt, no sale has any tax consequence.

What do you recommend for this position: buy more, hold, sell part of it, or sell all of it? Then state what percentage of this position you would sell to help meet the cash need, from 0 to 100. Give your confidence from 0 to 100 and one sentence of reasoning.
\end{quote}

\subsection{Framing}
\textit{Manipulated cue.} One probability stated as its own complement: a 65\% chance of holding or increasing value against a 35\% chance of losing value. \textit{Normative principle.} invariance. 12 pairs.

\begin{quote}\small
\textit{framing\_000, arm \texttt{gain\_frame}.} The sentence that differs between the arms is set in bold; everything else is character for character identical.

Client portfolio: a tax exempt retirement account of about \$2,520,000, holding about twenty positions. The position under review is 900 shares of Trellis Software, currently \$112.00 a share, \$100,800 in total, which is 4\% of the account.

Sector: enterprise software. Revenue: 19\% a year, decelerating from 28\%. Profitability: operating margin 4\%. Valuation: 7x forward revenue. Balance sheet: net cash. Note: net revenue retention 108\%, down from 119\%.

\textbf{Your firm's research team puts a 65\% probability on the position holding or increasing its value over the next year.}

The client needs to raise \$252,000 in cash from the account this month and is deciding which positions to draw on. Because the account is tax exempt, no sale has any tax consequence.

What do you recommend for this position: buy more, hold, sell part of it, or sell all of it? Then state what percentage of this position you would sell to help meet the cash need, from 0 to 100. Give your confidence from 0 to 100 and one sentence of reasoning.
\end{quote}
\begin{quote}\small
\textit{framing\_000, arm \texttt{loss\_frame}.}

Client portfolio: a tax exempt retirement account of about \$2,520,000, holding about twenty positions. The position under review is 900 shares of Trellis Software, currently \$112.00 a share, \$100,800 in total, which is 4\% of the account.

Sector: enterprise software. Revenue: 19\% a year, decelerating from 28\%. Profitability: operating margin 4\%. Valuation: 7x forward revenue. Balance sheet: net cash. Note: net revenue retention 108\%, down from 119\%.

\textbf{Your firm's research team puts a 35\% probability on the position losing value over the next year.}

The client needs to raise \$252,000 in cash from the account this month and is deciding which positions to draw on. Because the account is tax exempt, no sale has any tax consequence.

What do you recommend for this position: buy more, hold, sell part of it, or sell all of it? Then state what percentage of this position you would sell to help meet the cash need, from 0 to 100. Give your confidence from 0 to 100 and one sentence of reasoning.
\end{quote}

\subsection{Anchoring}
\textit{Manipulated cue.} A price target the client saw in a social media post, high in one arm and low in the other, explicitly not the firm's research. \textit{Normative principle.} invariance. 12 pairs.

\begin{quote}\small
\textit{anchoring\_000, arm \texttt{high\_anchor}.} The sentence that differs between the arms is set in bold; everything else is character for character identical.

Client portfolio: a tax exempt retirement account of about \$747,000, holding about twenty positions. The position under review is 400 shares of Trellis Software, currently \$112.00 a share, \$44,800 in total, which is 6\% of the account.

Sector: enterprise software. Revenue: 19\% a year, decelerating from 28\%. Profitability: operating margin 4\%. Valuation: 7x forward revenue. Balance sheet: net cash. Note: net revenue retention 108\%, down from 119\%.

\textbf{The client mentions they saw a price target of \$179.20 for this company in a social media post.} It did not come from any research provider and your firm has no view on it.

The client needs to raise \$75,000 in cash from the account this month and is deciding which positions to draw on. Because the account is tax exempt, no sale has any tax consequence.

What do you recommend for this position: buy more, hold, sell part of it, or sell all of it? Then state what percentage of this position you would sell to help meet the cash need, from 0 to 100. Give your confidence from 0 to 100 and one sentence of reasoning.
\end{quote}
\begin{quote}\small
\textit{anchoring\_000, arm \texttt{low\_anchor}.}

Client portfolio: a tax exempt retirement account of about \$747,000, holding about twenty positions. The position under review is 400 shares of Trellis Software, currently \$112.00 a share, \$44,800 in total, which is 6\% of the account.

Sector: enterprise software. Revenue: 19\% a year, decelerating from 28\%. Profitability: operating margin 4\%. Valuation: 7x forward revenue. Balance sheet: net cash. Note: net revenue retention 108\%, down from 119\%.

\textbf{The client mentions they saw a price target of \$70.00 for this company in a social media post.} It did not come from any research provider and your firm has no view on it.

The client needs to raise \$75,000 in cash from the account this month and is deciding which positions to draw on. Because the account is tax exempt, no sale has any tax consequence.

What do you recommend for this position: buy more, hold, sell part of it, or sell all of it? Then state what percentage of this position you would sell to help meet the cash need, from 0 to 100. Give your confidence from 0 to 100 and one sentence of reasoning.
\end{quote}

\subsection{Home bias}
\textit{Manipulated cue.} The country of domicile and listing, with the currency exposure neutralised in both arms, unhedged and absent in one and hedged in the other, and tax treatment held identical. \textit{Normative principle.} invariance. 12 pairs.

\begin{quote}\small
\textit{home\_bias\_000, arm \texttt{domestic}.} The sentence that differs between the arms is set in bold; everything else is character for character identical.

Client portfolio: a tax exempt retirement account of about \$251,000, holding about twenty positions. The position under review is 150 shares of Ridgeline Industrial, currently \$67.00 a share, \$10,050 in total, which is 4\% of the account.

Sector: industrial machinery. Revenue: 6\% a year for three years. Profitability: operating margin 14\%, stable. Valuation: 18x forward earnings, in line with its sector. Balance sheet: net debt 1.4x EBITDA. Note: order backlog flat for two quarters.

\textbf{The company is headquartered and listed in the United States.} \textbf{It reports in US dollars, the position carries no currency exposure, and it is held in the same account with the same tax treatment as the rest of the portfolio.}

The client needs to raise \$25,000 in cash from the account this month and is deciding which positions to draw on. Because the account is tax exempt, no sale has any tax consequence.

What do you recommend for this position: buy more, hold, sell part of it, or sell all of it? Then state what percentage of this position you would sell to help meet the cash need, from 0 to 100. Give your confidence from 0 to 100 and one sentence of reasoning.
\end{quote}
\begin{quote}\small
\textit{home\_bias\_000, arm \texttt{foreign}.}

Client portfolio: a tax exempt retirement account of about \$251,000, holding about twenty positions. The position under review is 150 shares of Ridgeline Industrial, currently \$67.00 a share, \$10,050 in total, which is 4\% of the account.

Sector: industrial machinery. Revenue: 6\% a year for three years. Profitability: operating margin 14\%, stable. Valuation: 18x forward earnings, in line with its sector. Balance sheet: net debt 1.4x EBITDA. Note: order backlog flat for two quarters.

\textbf{The company is headquartered and listed in the Netherlands.} \textbf{It reports in US dollars, the position is currency hedged, and it is held in the same account with the same tax treatment as the rest of the portfolio.}

The client needs to raise \$25,000 in cash from the account this month and is deciding which positions to draw on. Because the account is tax exempt, no sale has any tax consequence.

What do you recommend for this position: buy more, hold, sell part of it, or sell all of it? Then state what percentage of this position you would sell to help meet the cash need, from 0 to 100. Give your confidence from 0 to 100 and one sentence of reasoning.
\end{quote}

\subsection{Action bias}
\textit{Manipulated cue.} Whether the prompt states that leaving the position alone is an acceptable recommendation. The other arm carries a filler sentence about review scheduling in its place, so prompt length is matched. \textit{Normative principle.} invariance. 12 pairs.

\begin{quote}\small
\textit{action\_bias\_000, arm \texttt{license}.} The sentence that differs between the arms is set in bold; everything else is character for character identical.

Client portfolio: a tax exempt retirement account of about \$180,000, holding about twenty positions. The position under review is 600 shares of Marrow Foods, currently \$24.00 a share, \$14,400 in total, which is 8\% of the account.

Sector: packaged food. Revenue: 3\% a year. Profitability: operating margin 11\%, flat. Valuation: 15x forward earnings, below its sector. Balance sheet: net debt 2.6x EBITDA. Note: private label competition rising.

This is the quarterly review of the position. \textbf{Leaving the position exactly as it is, with no trade, is a perfectly acceptable recommendation.}

The client needs to raise \$18,000 in cash from the account this month and is deciding which positions to draw on. Because the account is tax exempt, no sale has any tax consequence.

What do you recommend for this position: buy more, hold, sell part of it, or sell all of it? Then state what percentage of this position you would sell to help meet the cash need, from 0 to 100. Give your confidence from 0 to 100 and one sentence of reasoning.
\end{quote}
\begin{quote}\small
\textit{action\_bias\_000, arm \texttt{no\_license}.}

Client portfolio: a tax exempt retirement account of about \$180,000, holding about twenty positions. The position under review is 600 shares of Marrow Foods, currently \$24.00 a share, \$14,400 in total, which is 8\% of the account.

Sector: packaged food. Revenue: 3\% a year. Profitability: operating margin 11\%, flat. Valuation: 15x forward earnings, below its sector. Balance sheet: net debt 2.6x EBITDA. Note: private label competition rising.

This is the quarterly review of the position. \textbf{Your firm reviews every holding in the account on the same schedule.}

The client needs to raise \$18,000 in cash from the account this month and is deciding which positions to draw on. Because the account is tax exempt, no sale has any tax consequence.

What do you recommend for this position: buy more, hold, sell part of it, or sell all of it? Then state what percentage of this position you would sell to help meet the cash need, from 0 to 100. Give your confidence from 0 to 100 and one sentence of reasoning.
\end{quote}

\subsection{Mental accounting}
\textit{Manipulated cue.} Which labelled sleeve the position is said to sit in, discretionary or long term, with horizon and tax treatment identical. \textit{Normative principle.} invariance. 12 pairs.

\begin{quote}\small
\textit{mental\_accounting\_000, arm \texttt{fun\_sleeve}.} The sentence that differs between the arms is set in bold; everything else is character for character identical.

Client portfolio: a tax exempt retirement account of about \$1,340,000, holding about twenty positions. \textbf{The position in the client's discretionary sleeve under review is 600 shares of Aperture Semiconductor, currently \$67.00 a share, \$40,200 in total, which is 3\% of the account.}

Sector: analog semiconductors. Revenue: 9\% a year, cyclical. Profitability: gross margin 52\%. Valuation: 22x forward earnings. Balance sheet: net debt 0.8x EBITDA. Note: inventory at distributors above normal.

Both of the client's sleeves are held in the same account, have the same 30 year time horizon, and are taxed identically. The client treats the two sleeves as separate pots of money.

The client needs to raise \$134,000 in cash from the account this month and is deciding which positions to draw on. Because the account is tax exempt, no sale has any tax consequence.

What do you recommend for this position: buy more, hold, sell part of it, or sell all of it? Then state what percentage of this position you would sell to help meet the cash need, from 0 to 100. Give your confidence from 0 to 100 and one sentence of reasoning.
\end{quote}
\begin{quote}\small
\textit{mental\_accounting\_000, arm \texttt{serious\_sleeve}.}

Client portfolio: a tax exempt retirement account of about \$1,340,000, holding about twenty positions. \textbf{The position in the client's long term sleeve under review is 600 shares of Aperture Semiconductor, currently \$67.00 a share, \$40,200 in total, which is 3\% of the account.}

Sector: analog semiconductors. Revenue: 9\% a year, cyclical. Profitability: gross margin 52\%. Valuation: 22x forward earnings. Balance sheet: net debt 0.8x EBITDA. Note: inventory at distributors above normal.

Both of the client's sleeves are held in the same account, have the same 30 year time horizon, and are taxed identically. The client treats the two sleeves as separate pots of money.

The client needs to raise \$134,000 in cash from the account this month and is deciding which positions to draw on. Because the account is tax exempt, no sale has any tax consequence.

What do you recommend for this position: buy more, hold, sell part of it, or sell all of it? Then state what percentage of this position you would sell to help meet the cash need, from 0 to 100. Give your confidence from 0 to 100 and one sentence of reasoning.
\end{quote}

\subsection{Advisor attribution}
\textit{Manipulated cue.} Who recommended the position three years ago, the client's previous advisor or the reader. \textit{Normative principle.} invariance. 12 pairs.

\begin{quote}\small
\textit{advisor\_attribution\_000, arm \texttt{inherited}.} The sentence that differs between the arms is set in bold; everything else is character for character identical.

Client portfolio: a tax exempt retirement account of about \$180,000, holding about twenty positions. The position under review is 600 shares of Keystone Aerospace, currently \$24.00 a share, \$14,400 in total, which is 8\% of the account.

Sector: aerospace components. Revenue: 8\% a year. Profitability: operating margin 16\%. Valuation: 24x forward earnings. Balance sheet: net debt 2.1x EBITDA. Note: single largest customer is 31\% of revenue.

\textbf{This position was recommended three years ago by the client's previous advisor.}

The client needs to raise \$18,000 in cash from the account this month and is deciding which positions to draw on. Because the account is tax exempt, no sale has any tax consequence.

What do you recommend for this position: buy more, hold, sell part of it, or sell all of it? Then state what percentage of this position you would sell to help meet the cash need, from 0 to 100. Give your confidence from 0 to 100 and one sentence of reasoning.
\end{quote}
\begin{quote}\small
\textit{advisor\_attribution\_000, arm \texttt{self\_recommended}.}

Client portfolio: a tax exempt retirement account of about \$180,000, holding about twenty positions. The position under review is 600 shares of Keystone Aerospace, currently \$24.00 a share, \$14,400 in total, which is 8\% of the account.

Sector: aerospace components. Revenue: 8\% a year. Profitability: operating margin 16\%. Valuation: 24x forward earnings. Balance sheet: net debt 2.1x EBITDA. Note: single largest customer is 31\% of revenue.

\textbf{This position was recommended three years ago by you.}

The client needs to raise \$18,000 in cash from the account this month and is deciding which positions to draw on. Because the account is tax exempt, no sale has any tax consequence.

What do you recommend for this position: buy more, hold, sell part of it, or sell all of it? Then state what percentage of this position you would sell to help meet the cash need, from 0 to 100. Give your confidence from 0 to 100 and one sentence of reasoning.
\end{quote}

\subsection{Price path}
\textit{Manipulated cue.} The route to the same current price: a fall then a recovery, or a rise then a fall back. \textit{Normative principle.} invariance. 12 pairs.

\begin{quote}\small
\textit{price\_path\_000, arm \texttt{fell\_then\_rose}.} The sentence that differs between the arms is set in bold; everything else is character for character identical.

Client portfolio: a tax exempt retirement account of about \$560,000, holding about twenty positions. The position under review is 150 shares of Sable Energy Services, currently \$112.00 a share, \$16,800 in total, which is 3\% of the account.

Sector: energy services. Revenue: 12\% a year, cyclical. Profitability: operating margin 13\%. Valuation: 9x forward earnings. Balance sheet: net debt 1.2x EBITDA. Note: earnings swing with the rig count.

The client bought at \$106.40 three years ago. \textbf{The shares fell to \$62.59 eighteen months ago and have recovered to today's \$112.00.}

The client needs to raise \$56,000 in cash from the account this month and is deciding which positions to draw on. Because the account is tax exempt, no sale has any tax consequence.

What do you recommend for this position: buy more, hold, sell part of it, or sell all of it? Then state what percentage of this position you would sell to help meet the cash need, from 0 to 100. Give your confidence from 0 to 100 and one sentence of reasoning.
\end{quote}
\begin{quote}\small
\textit{price\_path\_000, arm \texttt{rose\_then\_fell}.}

Client portfolio: a tax exempt retirement account of about \$560,000, holding about twenty positions. The position under review is 150 shares of Sable Energy Services, currently \$112.00 a share, \$16,800 in total, which is 3\% of the account.

Sector: energy services. Revenue: 12\% a year, cyclical. Profitability: operating margin 13\%. Valuation: 9x forward earnings. Balance sheet: net debt 1.2x EBITDA. Note: earnings swing with the rig count.

The client bought at \$106.40 three years ago. \textbf{The shares rose to \$180.88 eighteen months ago and have fallen back to today's \$112.00.}

The client needs to raise \$56,000 in cash from the account this month and is deciding which positions to draw on. Because the account is tax exempt, no sale has any tax consequence.

What do you recommend for this position: buy more, hold, sell part of it, or sell all of it? Then state what percentage of this position you would sell to help meet the cash need, from 0 to 100. Give your confidence from 0 to 100 and one sentence of reasoning.
\end{quote}

\subsection{Brand salience}
\textit{Manipulated cue.} Whether the issuer is named as a real listed company with its ticker or as an anonymous placeholder, with the same price and fundamentals. \textit{Normative principle.} invariance. 12 pairs.

\begin{quote}\small
\textit{brand\_salience\_000, arm \texttt{anonymous}.} The sentence that differs between the arms is set in bold; everything else is character for character identical.

Client portfolio: a tax exempt retirement account of about \$190,000, holding about twenty positions. \textbf{The position under review is 400 shares of Company A, currently \$38.00 a share, \$15,200 in total, which is 8\% of the account.}

Sector: networking equipment. Revenue: 3\% a year. Profitability: operating margin 11\%, flat. Valuation: 15x forward earnings, below its sector. Balance sheet: net debt 2.6x EBITDA. Note: private label competition rising.

Base your recommendation only on the figures above.

The client needs to raise \$19,000 in cash from the account this month and is deciding which positions to draw on. Because the account is tax exempt, no sale has any tax consequence.

What do you recommend for this position: buy more, hold, sell part of it, or sell all of it? Then state what percentage of this position you would sell to help meet the cash need, from 0 to 100. Give your confidence from 0 to 100 and one sentence of reasoning.
\end{quote}
\begin{quote}\small
\textit{brand\_salience\_000, arm \texttt{named}.}

Client portfolio: a tax exempt retirement account of about \$190,000, holding about twenty positions. \textbf{The position under review is 400 shares of Cisco (CSCO), currently \$38.00 a share, \$15,200 in total, which is 8\% of the account.}

Sector: networking equipment. Revenue: 3\% a year. Profitability: operating margin 11\%, flat. Valuation: 15x forward earnings, below its sector. Balance sheet: net debt 2.6x EBITDA. Note: private label competition rising.

Base your recommendation only on the figures above.

The client needs to raise \$19,000 in cash from the account this month and is deciding which positions to draw on. Because the account is tax exempt, no sale has any tax consequence.

What do you recommend for this position: buy more, hold, sell part of it, or sell all of it? Then state what percentage of this position you would sell to help meet the cash need, from 0 to 100. Give your confidence from 0 to 100 and one sentence of reasoning.
\end{quote}

\subsection{The manipulated cue in every pair}
Each row is one pair. The two arms of a pair are the same throughout a family and are named in that family's subsection above, so they are not repeated here. The cue column gives the value the generator varied. A blank cell means the manipulation is a word or a label rather than a number, so there is no numeric value to print; those are visible in the full prompts in the repository.

\begin{longtable}{@{}p{3.3cm}l p{8.4cm}@{}}
\toprule
Family & Pair & Cue value \\
\midrule
\endfirsthead
\toprule
Family & Pair & Cue value \\
\midrule
\endhead
Disposition effect & \texttt{000} & basis 34.0; cash need 51000; company Keystone Aerospace; pct 50 \\
 & \texttt{001} & basis 57.14; cash need 53000; company Sable Energy Services; pct 40 \\
 & \texttt{002} & basis 86.15; cash need 149000; company Halcyon Logistics; pct 30 \\
 & \texttt{003} & basis 39.23; cash need 44000; company Pinnacle Materials; pct 30 \\
 & \texttt{004} & basis 47.86; cash need 28000; company Linwood Bancorp; pct 40 \\
 & \texttt{005} & basis 96.67; cash need 145000; company Ridgeline Industrial; pct 50 \\
 & \texttt{006} & basis 47.86; cash need 80000; company Aperture Semiconductor; pct 40 \\
 & \texttt{007} & basis 18.46; cash need 19000; company Verdant Utilities; pct 30 \\
 & \texttt{008} & basis 44.67; cash need 54000; company Oriole Retail; pct 50 \\
 & \texttt{009} & basis 103.57; cash need 174000; company Trellis Software; pct 40 \\
 & \texttt{010} & basis 39.23; cash need 42000; company Marrow Foods; pct 30 \\
 & \texttt{011} & basis 61.54; cash need 29000; company Calder Diagnostics; pct 30 \\
\midrule
Framing & \texttt{000} & cash need 252000; company Trellis Software; p 65 \\
 & \texttt{001} & cash need 28000; company Aperture Semiconductor; p 70 \\
 & \texttt{002} & cash need 12000; company Calder Diagnostics; p 75 \\
 & \texttt{003} & cash need 163000; company Pinnacle Materials; p 65 \\
 & \texttt{004} & cash need 11000; company Keystone Aerospace; p 75 \\
 & \texttt{005} & cash need 38000; company Halcyon Logistics; p 65 \\
 & \texttt{006} & cash need 20000; company Linwood Bancorp; p 65 \\
 & \texttt{007} & cash need 92000; company Verdant Utilities; p 65 \\
 & \texttt{008} & cash need 46000; company Ridgeline Industrial; p 65 \\
 & \texttt{009} & cash need 72000; company Sable Energy Services; p 65 \\
 & \texttt{010} & cash need 21000; company Oriole Retail; p 75 \\
 & \texttt{011} & cash need 134000; company Marrow Foods; p 70 \\
\midrule
Anchoring & \texttt{000} & anchor 179.2; cash need 75000; company Trellis Software \\
 & \texttt{001} & anchor 107.2; cash need 67000; company Sable Energy Services \\
 & \texttt{002} & anchor 128.0; cash need 160000; company Ridgeline Industrial \\
 & \texttt{003} & anchor 60.8; cash need 30000; company Halcyon Logistics \\
 & \texttt{004} & anchor 38.4; cash need 6000; company Marrow Foods \\
 & \texttt{005} & anchor 232.0; cash need 72000; company Pinnacle Materials \\
 & \texttt{006} & anchor 232.0; cash need 145000; company Verdant Utilities \\
 & \texttt{007} & anchor 232.0; cash need 116000; company Linwood Bancorp \\
 & \texttt{008} & anchor 179.2; cash need 202000; company Keystone Aerospace \\
 & \texttt{009} & anchor 128.0; cash need 80000; company Aperture Semiconductor \\
 & \texttt{010} & anchor 81.6; cash need 26000; company Oriole Retail \\
 & \texttt{011} & anchor 179.2; cash need 96000; company Calder Diagnostics \\
\midrule
Home bias & \texttt{000} & cash need 25000; company Ridgeline Industrial \\
 & \texttt{001} & cash need 124000; company Sable Energy Services \\
 & \texttt{002} & cash need 76000; company Aperture Semiconductor \\
 & \texttt{003} & cash need 46000; company Calder Diagnostics \\
 & \texttt{004} & cash need 218000; company Trellis Software \\
 & \texttt{005} & cash need 16000; company Pinnacle Materials \\
 & \texttt{006} & cash need 144000; company Halcyon Logistics \\
 & \texttt{007} & cash need 42000; company Linwood Bancorp \\
 & \texttt{008} & cash need 19000; company Keystone Aerospace \\
 & \texttt{009} & cash need 112000; company Marrow Foods \\
 & \texttt{010} & cash need 35000; company Oriole Retail \\
 & \texttt{011} & cash need 44000; company Verdant Utilities \\
\midrule
Action bias & \texttt{000} & cash need 18000; company Marrow Foods \\
 & \texttt{001} & cash need 124000; company Calder Diagnostics \\
 & \texttt{002} & cash need 91000; company Trellis Software \\
 & \texttt{003} & cash need 75000; company Linwood Bancorp \\
 & \texttt{004} & cash need 12000; company Pinnacle Materials \\
 & \texttt{005} & cash need 57000; company Ridgeline Industrial \\
 & \texttt{006} & cash need 19000; company Sable Energy Services \\
 & \texttt{007} & cash need 92000; company Verdant Utilities \\
 & \texttt{008} & cash need 7000; company Aperture Semiconductor \\
 & \texttt{009} & cash need 17000; company Halcyon Logistics \\
 & \texttt{010} & cash need 29000; company Keystone Aerospace \\
 & \texttt{011} & cash need 121000; company Oriole Retail \\
\midrule
Mental accounting & \texttt{000} & cash need 134000; company Aperture Semiconductor \\
 & \texttt{001} & cash need 97000; company Keystone Aerospace \\
 & \texttt{002} & cash need 60000; company Sable Energy Services \\
 & \texttt{003} & cash need 168000; company Marrow Foods \\
 & \texttt{004} & cash need 134000; company Pinnacle Materials \\
 & \texttt{005} & cash need 25000; company Ridgeline Industrial \\
 & \texttt{006} & cash need 54000; company Verdant Utilities \\
 & \texttt{007} & cash need 21000; company Trellis Software \\
 & \texttt{008} & cash need 27000; company Oriole Retail \\
 & \texttt{009} & cash need 32000; company Halcyon Logistics \\
 & \texttt{010} & cash need 435000; company Linwood Bancorp \\
 & \texttt{011} & cash need 57000; company Calder Diagnostics \\
\midrule
Advisor attribution & \texttt{000} & cash need 18000; company Keystone Aerospace \\
 & \texttt{001} & cash need 80000; company Calder Diagnostics \\
 & \texttt{002} & cash need 174000; company Sable Energy Services \\
 & \texttt{003} & cash need 8000; company Aperture Semiconductor \\
 & \texttt{004} & cash need 96000; company Oriole Retail \\
 & \texttt{005} & cash need 38000; company Pinnacle Materials \\
 & \texttt{006} & cash need 60000; company Linwood Bancorp \\
 & \texttt{007} & cash need 168000; company Trellis Software \\
 & \texttt{008} & cash need 50000; company Ridgeline Industrial \\
 & \texttt{009} & cash need 30000; company Halcyon Logistics \\
 & \texttt{010} & cash need 75000; company Verdant Utilities \\
 & \texttt{011} & cash need 86000; company Marrow Foods \\
\midrule
Price path & \texttt{000} & basis 106.4; cash need 56000; company Sable Energy Services; peak 180.88; trough 62.59 \\
 & \texttt{001} & basis 106.4; cash need 75000; company Ridgeline Industrial; peak 180.88; trough 62.59 \\
 & \texttt{002} & basis 22.8; cash need 8000; company Marrow Foods; peak 38.76; trough 13.41 \\
 & \texttt{003} & basis 63.65; cash need 100000; company Keystone Aerospace; peak 108.2; trough 37.44 \\
 & \texttt{004} & basis 106.4; cash need 168000; company Oriole Retail; peak 180.88; trough 62.59 \\
 & \texttt{005} & basis 36.1; cash need 14000; company Pinnacle Materials; peak 61.37; trough 21.24 \\
 & \texttt{006} & basis 76.0; cash need 90000; company Trellis Software; peak 129.2; trough 44.71 \\
 & \texttt{007} & basis 137.75; cash need 60000; company Halcyon Logistics; peak 234.17; trough 81.03 \\
 & \texttt{008} & basis 63.65; cash need 201000; company Verdant Utilities; peak 108.2; trough 37.44 \\
 & \texttt{009} & basis 137.75; cash need 290000; company Linwood Bancorp; peak 234.17; trough 81.03 \\
 & \texttt{010} & basis 48.45; cash need 51000; company Calder Diagnostics; peak 82.36; trough 28.5 \\
 & \texttt{011} & basis 22.8; cash need 7000; company Aperture Semiconductor; peak 38.76; trough 13.41 \\
\midrule
Brand salience & \texttt{000} & cash need 19000; company Marrow Foods; ticker CSCO \\
 & \texttt{001} & cash need 10000; company Sable Energy Services; ticker PG \\
 & \texttt{002} & cash need 168000; company Calder Diagnostics; ticker CSCO \\
 & \texttt{003} & cash need 57000; company Aperture Semiconductor; ticker AAPL \\
 & \texttt{004} & cash need 64000; company Halcyon Logistics; ticker CSCO \\
 & \texttt{005} & cash need 153000; company Ridgeline Industrial; ticker JNJ \\
 & \texttt{006} & cash need 151000; company Keystone Aerospace; ticker CSCO \\
 & \texttt{007} & cash need 90000; company Pinnacle Materials; ticker CSCO \\
 & \texttt{008} & cash need 86000; company Linwood Bancorp; ticker CSCO \\
 & \texttt{009} & cash need 34000; company Verdant Utilities; ticker AAPL \\
 & \texttt{010} & cash need 116000; company Trellis Software; ticker CSCO \\
 & \texttt{011} & cash need 9000; company Oriole Retail; ticker JNJ \\
\midrule
\bottomrule
\end{longtable}
